% =============================================================================
% From Speakers to Dialects
% Cristina Gutierrez Alcala - Speech Signal Processing, 2026
%
% ACL style. Change the \usepackage line below to match the style file you
% already have in Overleaf (acl.sty, ACL2023.sty, emnlp2023.sty, ...).
% Upload the seven PNGs from output/figures/ into the Overleaf project root.
% =============================================================================

\documentclass[11pt]{article}

\usepackage[]{acl}

\usepackage{times}
\usepackage{latexsym}
\usepackage[T1]{fontenc}
\usepackage[utf8]{inputenc}
\usepackage{microtype}
\usepackage{graphicx}
\usepackage{booktabs}
\usepackage{makecell}
\usepackage{amsmath}
\usepackage{inconsolata}

\title{From Speakers to Dialects:\\An Acoustic Study of Northern and Southern
Peninsular Spanish Using Common Voice}

\author{Cristina Guti\'errez Alcal\'a \\
  Speech Signal Processing \\
  2026}

\begin{document}
\maketitle

\begin{abstract}
I compare Northern and Southern Peninsular Spanish acoustically, using the
dialect-labeled recordings of Mozilla Common Voice, with forced alignment for
phonemic measurement and all inference at the speaker level. Across 957 speakers
and 89,958 vowel tokens I find no difference between the two varieties in the
vowel space: nothing exceeds $|d| = 0.17$ in a design powered for
$d \geq 0.24$. A Mexican outgroup shows that this comes from the dialects and
not from the method, because the same procedure separates Mexican Spanish from
both Peninsular varieties with effects 2.5 times larger. The variation sits in
the consonant instead: Southern speech weakens /s/ in coda ($d = +0.63$), with
onset /s/ as a null control. I also document how an earlier version of this
analysis produced effects of more than one standard deviation that came from two
individual speakers, and what a volunteer corpus requires if that error is not to
be repeated.
\end{abstract}

% =============================================================================
\section{Introduction}

Every Spanish speaker can hear that someone from Bilbao and someone from Seville
sound different, and dialectology has described that difference for over a
century. Turning it into a reproducible acoustic measure is harder, and that is
the problem this paper works on.

The first two formants of a vowel, F1 and F2, are the resonance frequencies of
the vocal tract that encode its height and frontness, and they are the standard
way of representing a variety's vowel space. My initial question was whether the
vowel space of Northern peninsular Spanish differs measurably from that of
Southern peninsular Spanish, and whether that difference is enough to classify an
unknown speaker. Mozilla Common Voice looked like the right resource: thousands
of openly available recordings with dialect labels declared by the speakers
themselves, on a scale that laboratory phonetics rarely reaches.

A first version of the analysis said yes, and said it loudly: differences of more
than one and a half standard deviations in F1 and in the bandwidth of F2, with
$p$-values below 0.001. An audit then showed that those effects came from two
specific speakers, who between them accounted for 54\% and 80\% of their
respective dialect cells. The study did not compare two varieties; it compared
two people.

What follows is the analysis rebuilt around that discovery, with the discovery
itself treated as part of the subject rather than as something to hide. I report
four things. Section~\ref{sec:pseudo} documents how clip-based sampling produced
the large spurious effects, and shows that their magnitude depends directly on
how much one individual is allowed to dominate the sample.
Section~\ref{sec:vowelspace} uses a balanced design of 957 speakers and finds no
difference at all between the two peninsular varieties in the vowel space, in a
contrast powerful enough to rule out anything beyond a small effect.
Section~\ref{sec:outgroup} checks that this absence is not just a blunt method,
by adding Mexican Spanish as an outgroup that the same procedure does separate
from both varieties. Section~\ref{sec:coda} then finds a substantial effect in
the consonant system: the weakening of /s/ in coda in Southern speech.

What links the four is design rather than statistics, and I return to it in
Section~\ref{sec:discussion}. In a corpus built from very unequal contributions
by volunteers, contrasts drawn inside a single speaker hold up, and overall
comparisons between groups do not.

% =============================================================================
\section{State of the art}

\subsection{Dialectal variations in Peninsular Spanish}
\label{sec:dialectal}

Spanish has a very stable system of five vowels: /a, e, i, o, u/. It has no
distinctions of vowel length or tension, its unstressed vowels are not reduced as
they are in English or Portuguese, and its inventory has barely changed since the
Romance system became fixed. That stability matters here, because it leaves the
vowel system with little room for dialectal differentiation: variation in Spanish
is overwhelmingly consonantal \citep{hualde2005,martinez2007}. The reference
formant values for peninsular Spanish come from \citet{quilis1981}, measured on
laboratory speech, and are still the standard basis for comparison.

The best-documented feature of Southern speech is the weakening of /s/ in
implosive position, realized as aspiration, assimilation to the following
segment, or elision, against the retention that characterizes the peninsular
North \citep{alvar1996,narbona1998,gerfen2002,ruch2014}. A classic descriptive
account ties this weakening to the preceding vowel, which is said to compensate
for the loss by opening and lengthening, to the point of positing an expanded
vowel system based on degree of openness for Eastern Andalusian
\citep{alonso1950}. That account rests on auditory observation and small samples,
and current methods can test it directly (Section~\ref{sec:precoda}).

\subsection{Formant Normalization}
\label{sec:norm}

Formant frequencies depend on the length of the vocal tract as much as on
articulatory configuration, so two speakers producing the same vowel differ
systematically in hertz if they differ in anatomy. Comparing varieties in raw
hertz, on data that are not balanced by sex, confuses dialect with physiology.
Sociophonetics deals with this by normalizing per speaker, and of the available
procedures the intra-speaker z-score of \citet{lobanov1971} performs best in
systematic comparisons: it removes anatomical variation and keeps sociolinguistic
differences \citep{adank2004}. I use the perceptually based Bark scale
\citep{traunmuller1990} as a complementary measure.

\subsection{Common Voice as a phonetic corpus}
\label{sec:corpus}

Common Voice is a multilingual corpus of read speech built from volunteer
contributions and distributed openly \citep{ardila2020}. For dialectology it has
something rare: variety labels declared by the speakers themselves, across
thousands of speakers. It also has two problems. The labels are neither
verifiable nor granular, and, more seriously, contributions are spread very
unevenly, because nothing in the project's design pushes towards an even
distribution: a few enthusiastic volunteers can each record hundreds of clips.
Sampling that picks clips without grouping them by speaker therefore produces
observations that are not independent, which risks pseudoreplication
\citep{hurlbert1984,winter2019}. Good forced aligners for Spanish
\citep{mcauliffe2017} are what make it possible to work with phoneme-level
measures instead of clip-level ones, and so to look at the intra-speaker
contrasts that this corpus does support.

% =============================================================================
\section{Corpus and sampling}

\subsection{Common Voice 26.0 and dialect tags}

All the data used come from Mozilla Common Voice, Spanish 26.0 version (12 of
June of 2026), a corpus of read speech built from voluntary submissions and
distributed under an open licence \citep{ardila2020}. The version used contains
439,225 validated clips. Each clip has an orthographic transcription associated
with it and a speaker identification number (\texttt{client\_id}), in addition to
autodeclared metadata such as gender, age and accent.

Getting the corpus onto a machine where I could work with it shaped the pipeline
more than I expected. The Spanish release is a single archive of roughly 48 GB.
My first attempt was to download it, unpack it and work from the audio. That cost
two days of failed transfers and still left me without room for the extracted
files. The fix turned into the design. Nothing in the sampling stage needs the
audio, because the selection is made entirely from \texttt{validated.tsv}, which
sits at the front of the archive and can be pulled out with
\texttt{tar -{}-fast-read} in about a second. Only the clips the sampler selects
are then extracted, which in the final study meant unpacking 5,577 files instead
of 439,225.

The accent tag is what makes this corpus usable for a dialectal study, since
speakers tag the variety they consider themselves to speak. From the available
Spanish tags I took three: ``peninsular septentrional'', which I call North,
``peninsular meridional'', which I call South, and Mexican, which I include as an
external group for the reasons given in Section~\ref{sec:outgroup}. Speakers can
also mark more than one accent, and a fair number do. I discarded those clips
rather than assigning one tag by hand, because any tie-breaking rule would put my
own judgement into the grouping variable. That left 236,832 clips carrying
exactly one of the three tags.

These tags are self-declarations, not dialectological classifications. Nobody
verifies where a speaker is from, and nothing finer than the region is recorded.
I come back to this limitation in Section~\ref{sec:limitations}.

\subsection{Speaker-level sampling and per-speaker cap}

The most important design choice in this project was also the simplest: I sample
speakers, not clips. An older version of the analysis picked clips at random,
identifying them by file path, without looking at who had recorded each one. The
result was a sample dominated by a handful of individuals, and
Section~\ref{sec:pseudo} deals with the consequences.

The current sampling groups first by \texttt{client\_id}, orders each speaker's
clips by a reproducible criterion, and caps them at eight per speaker. The aim is
to maximize the number of speakers first and the number of clips per speaker
second, because the speaker count, not the token count, is what bounds the
inference. The whole stage is seeded and recorded in a versioned manifest, so the
sample can be rebuilt clip by clip.

The cap of eight is a compromise between two requirements that pull in opposite
directions. Lobanov normalization needs enough vowels from each speaker to
estimate that speaker's own mean and spread, which argues for more clips per
person. Independence argues for fewer. Eight clips of roughly five seconds yield
on the order of a hundred vowel tokens per speaker, comfortably enough for the
normalization, while holding any single individual's share of a dialect cell
below one percent. The sampling is implemented with pandas \citep{mckinney2010}
over NumPy arrays \citep{harris2020}.

One consequence of prioritizing speakers is that the design is bounded by its
smallest group. Southern Peninsular Spanish has 188 speakers in release 26.0, and
that number constrains everything the study can claim, so Northern and Southern
speakers were both taken in full. Mexican Spanish has over 1,500 available, far
more than the comparison needs, and was capped at 300.

I measure the effect of the cap with the Gini coefficient of clips per speaker,
where 0 means a perfectly even distribution and 1 means total concentration. In
the preliminary sample it was 0.76 in North and 0.85 in South. In the current
sample it is 0.26 and 0.18. Everything that follows depends on that change.

\subsection{Final design of the sample}
\label{sec:sample}

The resulting sample has 5,577 clips of 957 different speakers, and a total of
484.5 recorded minutes. Table~\ref{tab:sample} compiles its composition.

\begin{table}[t]
\centering
\small
\setlength{\tabcolsep}{4pt}
\begin{tabular}{lrrrrr}
\toprule
Variety & Clips & Speakers & \makecell{With\\tokens} & \makecell{Norma-\\lizable} & Tokens \\
\midrule
North  & 2,639 & 469 & 468 & 443 & 42,805 \\
South  & 1,199 & 188 & 188 & 181 & 19,549 \\
Mexico & 1,739 & 300 & 299 & 282 & 27,604 \\
\midrule
Total  & 5,577 & 957 & 955 & 906 & 89,958 \\
\bottomrule
\end{tabular}
\caption{Composition of the sample in each variety.}
\label{tab:sample}
\end{table}

The imbalance between varieties is dictated by the corpus itself. The Southern
part of the peninsula is represented by fewer than half as many speakers as the
North. I do not correct for this by undersampling the North, because discarding
valid data solely to balance the cells would reduce statistical power without
gaining anything. The comparisons are calculated using tests that do not assume
equal sample sizes.

The imbalance by sex does matter. Male voices dominate all three varieties: 27.9\%
of speakers are female in the North, 20.4\% in the South, and 25.2\% in Mexico,
with the rest male or undeclared. Since vocal tract length influences formant
frequencies directly, a comparison between these groups in raw hertz cannot
separate dialect from anatomy. The normalization described in
Section~\ref{sec:qc} is therefore a requirement of the design rather than a
refinement.

Figure~\ref{fig:dataflow} summarizes the complete path, from the 439,225
validated clips in the corpus to the tokens finally analyzed. The alignment and
quality-control stages that produce those numbers are described in
Sections~\ref{sec:align} and~\ref{sec:qc}.

\begin{figure*}[t]
\centering
\includegraphics[width=\textwidth]{fig_data_flow.png}
\caption{From corpus to analyzed tokens. The left panel shows the reduction from
the validated Common Voice clips to the final sample after filtering by unique
label and applying the limit of eight clips per speaker. The right panel shows
the quality control funnel applied to the extracted vowel tokens.}
\label{fig:dataflow}
\end{figure*}

The sample of 903 clips on which the preliminary analysis was based remains
unchanged and is reused in Section~\ref{sec:pseudo}, where it serves not as a
data source but as an object of study.

% =============================================================================
\section{Methodology}

\subsection{Audio preprocessing and forced alignment}
\label{sec:align}

Each clip goes from MP3 to mono WAV at 16 kHz, the sampling frequency the aligner
expects, with peak scaling so that resampling does not clip the signal. Beside
each audio file I write the transcription of what the speaker read aloud, taken
straight from the corpus metadata, and I lay the pairs out in one directory per
speaker so that the aligner can apply its speaker adaptation.

Forced alignment was not part of the original plan. It became necessary after two
attempts to recover vowel identity by cheaper means, and both failures explain
why the pipeline looks the way it does. The first classified each measurement by
its position in the F1/F2 plane, using the reference ranges for Spanish vowels in
\citet{quilis1981}. Nearly every token it labelled /a/ came from the Northern
group and nearly every /i/ from the Southern one. Those ranges are calibrated on
Northern speech, so a displaced Southern token lands in a neighbouring category.
The procedure is also circular: labelling vowels by their formants and
then comparing formants across dialects builds the conclusion into the labels.
The second attempt inferred the vowel from the transcription, which works only
for clips containing a single vowel type. Of 903 clips, 14 resolved.

That left aligning the audio to its transcription properly. I installed the
Montreal Forced Aligner \citep{mcauliffe2017}, built on Kaldi \citep{povey2011},
together with the pretrained \texttt{spanish\_mfa} acoustic model and dictionary.
The aligner takes the audio and its orthographic transcription and returns a
TextGrid with the time boundaries of each word and each phone. Of the 5,577
clips, 5,572 aligned correctly (99.9\%); the remaining five failed because of
discrepancies between the transcription and the signal.

With alignment, each clip yields all of its vowel nuclei with phonemic labels,
and with them four pieces of information that were unavailable before: vowel
identity, duration, stress position, and the following segment. The last of these
is what made the coda /s/ analysis of Section~\ref{sec:coda} possible. The
semivowels /j/ and /w/ are deliberately excluded, since they are not syllable
nuclei and their formant targets are transitional.

\subsection{Extraction of acoustic parameters}

The acoustic measures were taken using Praat \citep{boersma2021} through its
Python interface Parselmouth \citep{jadoul2018}, using linear predictive analysis
with Burg's method with a maximum of five formants.

The frequency ceiling follows the speaker's declared sex: 5,000 Hz for male
voices, 5,500 Hz for female voices, and 5,250 Hz when the corpus records no
value. The previous analysis used a single ceiling of 5,000 Hz for everyone,
which biased the estimates for female voices, and that matters here because the
dialect cells are not balanced by sex.

From each vowel I take F1, F2, B1 and B2 at three points: 25\%, 50\% and 75\% of
its duration. The midpoint is the conventional measure, and the other two let me
quantify how much the formants move inside the vowel, which is relevant to
descriptions of diphthongization in Southern speech. To these I add segment
duration, fundamental frequency, the harmonics-to-noise ratio (HNR), the clip's
signal-to-noise ratio, lexical stress and the following consonantal context.
Vowels shorter than 30 ms are dropped at the start, since they consist of
transitions into the adjacent consonants and have no steady state worth
measuring. The extraction runs in parallel across processes.

\subsection{Quality control and normalization}
\label{sec:qc}

Quality control first applies deliberately generous range filters: F1 between 150
and 1,100 Hz, F2 between 500 and 3,200 Hz, and bandwidth between 10 and 900 Hz
(B1) and 10 and 1,000 Hz (B2). These limits exist to remove tracking failures (an
F1 of 3 kHz is an octave error, not a vowel), not to impose a phonetic
expectation. Tightening them further would reintroduce the circularity that
invalidated the vowel tagging of the preliminary analysis. Also, F2 must be
higher than F1 and outliers beyond three standard deviations within each speaker
and each vowel were eliminated, a filter that is only possible when the identity
of the segment is known.

\begin{table}[t]
\centering
\small
\begin{tabular}{lrr}
\toprule
Phase & Tokens & \% retained \\
\midrule
Extracted                        & 100,053 & 100.0 \\
Measurable at midpoint           & 100,052 & 100.0 \\
F1/F2 in range                   &  99,104 &  99.1 \\
B1/B2 in range                   &  91,198 &  91.1 \\
F2 $>$ F1                        &  91,198 &  91.1 \\
Within-speaker outliers removed  &  89,958 &  89.9 \\
\bottomrule
\end{tabular}
\caption{Quality control funnel on the extracted vowel tokens.}
\label{tab:qc}
\end{table}

Most of the loss comes from the bandwidth filters, which discard 8\% of the
tokens. Bandwidths estimated by linear prediction turned out to be unstable and
sensitive to recording conditions, which matters in Section~\ref{sec:pseudo},
where one of the effects of the preliminary analysis was an effect of B2.

On the surviving tokens, three representations were calculated. The first is raw
hertz, which is retained for transparency but does not support any comparisons
between groups, for the reasons explained in Section~\ref{sec:sample}. The second
is the Bark scale \citep{traunmuller1990}, which approximates the ear's frequency
resolution. The third, and the one used in all contrasts, is Lobanov
normalization \citep{lobanov1971}: the z-score of each formant computed within
each speaker across all of their vowels. It is the best-performing procedure in
systematic comparisons of normalization methods \citep{adank2004}, for the
reasons set out in Section~\ref{sec:norm}. It requires that the speaker has
produced all five vowel categories, a condition met by 906 of the 955 speakers
with tokens. The remaining 49 are excluded from the normalized analyses.

\subsection{Statistical scope}

All the inference here rests on one principle: the unit of analysis is the
speaker. Tokens from the same speaker are not independent observations, since
they share an anatomy, a microphone, a room and a reading style. Treating them as
independent inflates the degrees of freedom and produces significance whether or
not a dialectal effect exists, as Section~\ref{sec:pseudo} shows empirically.

In practice, this means that the data are aggregated into speaker $\times$ vowel
cells and, from there, into a single measure per speaker, prior to any
comparisons. Comparisons between varieties are performed using Welch's $t$-test,
which does not assume equal variances or sample sizes, along with Cohen's $d$
\citep{cohen1988} and confidence intervals obtained via bootstrapping with 2,000
resamples of speakers, not tokens. Welch's test and the bootstrap are computed
with SciPy \citep{virtanen2020}, and the mixed models are fitted with statsmodels
\citep{seabold2010}. Within each family of contrasts, $p$-values are corrected
for multiple comparisons using Holm's procedure \citep{holm1979}. Mixed models
with a random intercept per speaker are fitted to the means per speaker and
vowel. The significance level is 0.05, and every stochastic step, like sampling,
bootstrapping, $k$-means initialization, and cross-validation partitions, is
seeded with a fixed seed.

One methodological trap deserves a mention, because I nearly fell into it.
Fitting a token-level mixed model to values already normalized by Lobanov yields
an ICC close to zero and $p$-values below 0.001. Lobanov sets each
speaker's mean to 0 and their standard deviation to 1, so that there is no
remaining variance between speakers for a random intercept to account for, and
the model reverts to treating the tokens as independent. A random intercept
protects against speakers differing in level, not against one speaker shaping a
within-speaker pattern. Aggregating into speaker $\times$ vowel cells before
modeling is what solves the problem.

Finally, I compute a power analysis from the speaker counts actually obtained,
which sets a limit on what the study can honestly claim. The thresholds appear
with the results in Section~\ref{sec:vowelspace}.

\subsection{Validation of the extractor}

To check that the extractor measures what it claims to measure, I synthesize
vowels with known formants using an F1--F5 resonator bank on a 120 Hz glottal
source, and compare the recovered value against the real one. It is accurate for
open vowels, where the error in F1 for /a/ is +0.2\%, and degrades for closed
ones, with errors of +13.5\% for /i/ and +33.7\% for /u/. The cause is well
known: in closed vowels F1 approaches F0, so individual harmonics dominate the
fitted envelope, and in /u/ the F1 and F2 poles interfere with each other.

The bias is identical in both dialects, so it cannot generate a spurious
difference between them, but it makes /u/ the least reliable measure in the
study, and it is treated as such in the discussion. The test is part of a battery
of 24 automated tests that also verify the handling of coda context, the
identification of the stressed syllable, and the normalization functions.

The pipeline runs as a single seeded entry point; figures are produced with
Matplotlib \citep{hunter2007}, and TextGrids are parsed with praatio. A full run
takes about fifteen minutes end to end.

% =============================================================================
\section{Pseudoreplication}
\label{sec:pseudo}

\subsection{Naive analysis}

The first version of this study worked on 440 tokens (244 North and 196 South),
one measurement per clip taken at the global intensity peak. A Welch $t$-test on
those tokens made three of the four comparisons significant, two of them with
effects of a size that dialectal phonetics rarely sees: F1 at $d = +1.34$ and B2
at $d = -1.43$, both with $p < 0.001$. Read at face value, those numbers describe
two varieties whose vowels sit more than one and a half standard deviations
apart.

The result looked excellent. It was also false, and the way it was false turned
out to be more useful than any effect it could have found.

\subsection{Audit per speaker}

The 440 tokens do not come from 440 different speakers. Cross-referencing the
measurements with the speaker identifiers in the corpus gives 54 Northern
speakers and 26 Southern ones, distributed very unevenly. One speaker contributed
132 tokens, 54\% of the Northern cell, and another 157, 80\% of the Southern one.
The three most prolific speakers in each cell hold 64\% and 86\% of the data.

Removing speakers one at a time shows what this does. Drop the dominant Southern
speaker and $d$(F1) falls from 1.34 to 0.23, $d$(B2) from $-1.43$ to $-0.45$.
Drop the Northern one and they fall to 1.04 and $-1.04$. Drop both, which means
289 of the 440 tokens and leaves 151, and the effects come down to $d = +0.01$
and $d = -0.09$, which is nothing. Removing any other speaker leaves the numbers
almost untouched.

What the preliminary analysis measured, then, was the difference between two
particular people rather than between two varieties. \citet{hurlbert1984} called
this pseudoreplication: treating observations that share an experimental unit as
independent replicates. The $t$-test assumes independence, 440 clips from 80
speakers do not provide it, and when one person supplies 80\% of a cell the
effective size of that cell is closer to one than to two hundred.

\subsection{Size of the effect}

The dependence can be directly quantified by recalculating the original analysis
under increasing token caps per speaker. With a cap of one token per speaker (80
strictly independent observations) $d$(F1) is $+0.11$. With a cap of three,
$+0.11$. With a cap of eight, $+0.30$. With a cap of twenty, $+0.54$; with a cap
of fifty, $+0.90$; with a cap of one hundred, $+1.15$; and with no cap at all,
$+1.34$. The relationship is monotonic, and it is replicated in B2 with the
opposite sign.

The effect size I reported was therefore a direct function of how much one
individual was allowed to dominate the sample. There is no ``true'' value that
the cap distorts. The cap decides the value.

\subsection{Naive against corrected}

Table~\ref{tab:naive} contrasts the four original comparisons with two analyses
that respect the data structure: the aggregation of one measure per speaker and a
mixed model with a random intercept per speaker \citep{barr2013,winter2019}.

\begin{table}[t]
\centering
\small
\setlength{\tabcolsep}{4pt}
\begin{tabular}{lrrrrrr}
\toprule
 & \makecell{naive\\$d$} & \makecell{$p$\\(Holm)} &
 \makecell{per spk\\$d$} & $p$ & \makecell{mixed\\$p$} & ICC \\
\midrule
F1 & $+1.335$ & $<$0.001 & $+0.040$ & 1.000 & 1.000 & 0.54 \\
F2 & $-0.337$ & 0.001    & $+0.333$ & 0.835 & 1.000 & 0.17 \\
B1 & $-0.166$ & 0.100    & $+0.195$ & 1.000 & 1.000 & 0.05 \\
B2 & $-1.433$ & $<$0.001 & $-0.172$ & 1.000 & 1.000 & 0.28 \\
\bottomrule
\end{tabular}
\caption{Naive analysis per token against a speaker level analysis, on 440
tokens.}
\label{tab:naive}
\end{table}

No effect survives. F2 also changes sign at the speaker level, on top of losing
its significance, so it cannot be read as a real effect that has merely been
attenuated.

The ICC for F1 is 0.54: more than half of the variance in F1 comes from who is
speaking rather than from which dialect they speak. The design effect computed
from that value says the 440 pooled tokens carry the information of roughly 128
independent observations. The bootstrap intervals show the same thing from
another angle. When resampling tokens,
$d$(F1) has an interval of $[+1.11, +1.60]$, which comfortably excludes zero.
When resampling speakers, the interval is $[-0.20, +1.83]$ and includes it. For
B2, $[-1.67, -1.22]$ versus $[-1.89, +0.19]$. All four speaker-resampled
intervals contain zero.

\begin{figure*}[t]
\centering
\includegraphics[width=0.75\textwidth]{fig_pseudoreplication.png}
\caption{Anatomy of a sampling artifact, based on the 440 tokens from the
preliminary analysis. A: A few speakers contribute most of the data for each
dialect. B: The effects fade as the cap on tokens per speaker prevents any single
individual from dominating. C: Removing a single prolific speaker eliminated the
effects. D: The confidence intervals obtained by resampling speakers all include
zero, unlike those obtained by resampling tokens.}
\label{fig:pseudo}
\end{figure*}

\subsection{Why the balanced result helps}

When applied to the sample described in Section~\ref{sec:sample}, the same audit
procedure finds nothing to audit. The most prolific southern speaker now accounts
for 1\% of their cell, and their removal shifts the largest per-vowel effect by
0.00, compared to the shift from 1.6 to 0.0 produced by the same operation in the
previous sample. The pooled analysis and the speaker-level analysis coincide,
simply because there is nothing left on which they might disagree.

This convergence is the argument in favor of the per-speaker cap, and it is also
the reason why the results in Section~\ref{sec:results} can be interpreted with
confidence: not because the statistical method is more sophisticated, but because
the sampling design no longer allows the question to answer itself.

% =============================================================================
\section{Results}
\label{sec:results}

\subsection{Vowel space on speaker level}
\label{sec:vowelspace}

The main comparison runs on one average per speaker and per measure, over 456 to
461 Northern speakers and 178 to 186 Southern ones, depending on how many
produced the vowel in question. Table~\ref{tab:speakerlevel} gathers the fifteen
measures: the normalized positions of the five vowels in F1 and F2, mean
duration, the bandwidths, the harmonics-to-noise ratio and the fundamental
frequency.

\begin{table}[t]
\centering
\small
\setlength{\tabcolsep}{4pt}
\begin{tabular}{lrrrrr}
\toprule
Measurement & North & South & $d$ & $p$ & \makecell{$p$\\(Holm)} \\
\midrule
F1 /a/       & 1.09   & 1.07   & $+0.12$ & 0.166 & 1.000 \\
F1 /e/       & $-0.31$ & $-0.28$ & $-0.17$ & 0.043 & 0.638 \\
F1 /i/       & $-1.22$ & $-1.21$ & $-0.03$ & 0.720 & 1.000 \\
F1 /o/       & $-0.09$ & $-0.08$ & $-0.04$ & 0.661 & 1.000 \\
F1 /u/       & $-0.84$ & $-0.84$ & $+0.01$ & 0.893 & 1.000 \\
F2 /a/       & $-0.25$ & $-0.27$ & $+0.14$ & 0.085 & 1.000 \\
F2 /e/       & 0.59   & 0.58   & $+0.02$ & 0.783 & 1.000 \\
F2 /i/       & 1.50   & 1.52   & $-0.06$ & 0.440 & 1.000 \\
F2 /o/       & $-0.96$ & $-0.95$ & $-0.04$ & 0.684 & 1.000 \\
F2 /u/       & $-1.10$ & $-1.16$ & $+0.11$ & 0.135 & 1.000 \\
Duration (s) & 0.077  & 0.077  & $-0.00$ & 0.995 & 1.000 \\
B1 (Hz)      & 149.2  & 146.1  & $+0.05$ & 0.528 & 1.000 \\
B2 (Hz)      & 218.0  & 216.8  & $+0.02$ & 0.822 & 1.000 \\
HNR (dB)     & 9.97   & 9.81   & $+0.04$ & 0.616 & 1.000 \\
F0 (Hz)      & 151.3  & 149.4  & $+0.05$ & 0.580 & 1.000 \\
\bottomrule
\end{tabular}
\caption{Speaker level comparison between North and South. F1 and F2 values are
normalized by Lobanov (intra-speaker z-scores); the bandwidths, the duration, the
HNR and F0 are expressed using their original units. A positive value of $d$
indicates a higher value in the North.}
\label{tab:speakerlevel}
\end{table}

None of these comparisons reaches significance once corrected for multiple
comparisons. The largest effect anywhere in the table is $|d| = 0.17$, for F1 of
/e/, whose uncorrected $p$ of 0.043 becomes 0.638 after Holm over the family of
fifteen contrasts. The bandwidths, where the preliminary analysis had put its
second strongest effect ($d = -1.43$), now give $d = +0.02$.

Global measures of the vowel space give the same answer as the individual
positions. The area of the convex hull defined by the five normalized vowels is
3.05 in the North and 2.95 in the South, with $d = +0.12$, a 95\% confidence
interval of $[-0.05, +0.28]$ and $p = 0.167$. Mixed models fitted to the speaker
$\times$ vowel means show no significant main effect of dialect.

\begin{figure*}[t]
\centering
\includegraphics[width=\textwidth]{fig_vowel_space.png}
\caption{Vowel space of the Northern and Southern parts of the Iberian Peninsula.
On the left, in raw hertz; on the right, after Lobanov normalization. Each point
represents the average for a single speaker, and the ellipses cover 95\% of the
variation among speakers. The reduction in dispersion between the two panels
corresponds to the anatomical variation that normalization eliminates; what
remains is nearly a complete overlap between varieties.}
\label{fig:vowelspace}
\end{figure*}

The power of the design is what makes this result informative, more than the
absence of significance itself. With 461 Northern and 186 Southern speakers, the
comparison would catch any effect of $d \geq 0.24$ with 80\% probability, and any
effect of $d \geq 0.28$ with 90\%. The preliminary sample could only resolve
effects of $d \geq 0.63$, which is to say large ones.

The conclusion has to be stated precisely. The two vowel spaces are not
identical: differences below $d \approx 0.24$ remain entirely possible and this
study cannot rule them out. What I can state is that no difference between the
North and the South of the Peninsula is large enough for a design with more than
six hundred speakers to detect. Whatever effect exists is small at most.

That leaves the obvious objection, which is that the method simply cannot tell
dialects apart.

\subsection{Calibration with an external group}
\label{sec:outgroup}

The result in Section~\ref{sec:vowelspace} allows for two incompatible
interpretations. One is that the two peninsular varieties do not differ
appreciably in their vowel systems. The other is that the procedure is too crude
to detect any dialectal differences, and would have produced a null result
regardless of the comparison made. Without additional evidence, both explanations
account for the data equally well.

The way to distinguish between them is to apply exactly the same procedure to a
contrast whose existence is not in dispute. To that end, 300 Mexican speakers
were included in the design and processed using the same sampling, the same
alignment, the same feature extraction, the same normalization, and the same
statistical tests. If the instrument distinguishes between Peninsular and Mexican
Spanish but not between Northern and Southern Spanish, the absence of an effect
in Section~\ref{sec:vowelspace} is a statement about the dialects, not about the
method.

The three pairwise contrasts, based on the twelve vowel measures evaluated at the
speaker level, show the following results. North versus South gives an average
$|d|$ of 0.064 and a maximum of 0.172, with 0 of the 12 comparisons significant.
North versus Mexico gives 0.152 and 0.303, with 4 of 12. South versus Mexico
gives 0.173 and 0.411, with 2 of 12.

Transatlantic contrasts produce effects that are, on average, 2.5 times greater
than those of the Peninsular contrast, and reach statistical significance after
correction in six of the twenty-four comparisons, whereas the North-South
contrast does not reach significance in any of the twelve. The largest
differences are concentrated, as might be expected, in F2 of the peripheral
vowels: F2 of /a/ between the South and Mexico ($d = -0.41$, $p < 0.001$), F2 of
/i/ between the South and Mexico ($d = +0.40$, $p < 0.001$), and F2 of /i/
between the North and Mexico ($d = +0.30$, $p = 0.004$). Vowel duration also
distinguishes the North from Mexico ($d = -0.28$, $p = 0.032$), without
distinguishing the North from the South.

\begin{figure*}[t]
\centering
\includegraphics[width=\textwidth]{fig_outgroup.png}
\caption{Calibration with an external group. The three varieties overlaid in the
normalized vowel space and the average magnitude of the effect for each pairwise
contrast. The peninsular contrast is consistently the smallest of the three.}
\label{fig:outgroup}
\end{figure*}

The conclusion is that the method does resolve dialectal differences when they
exist, using the same sample, the same measures, and the same decision threshold.
The null result in Section~\ref{sec:vowelspace} is not an artifact of a blunt
instrument.

These effects are small in absolute terms, and this is the place to say so rather
than the discussion. A mean $|d|$ of 0.16 for the transatlantic contrast is not
what the descriptive literature would lead you to expect from two varieties
separated by an ocean and five centuries. I read that as a problem of scale
rather than a failure of the comparison. The five-vowel system described in
Section~\ref{sec:dialectal} gives dialectal drift very little to work with, and
the outgroup is what makes this visible: even across the Atlantic, these vowels
barely move. Seen that way, the null result of Section~\ref{sec:vowelspace} is
close to what Spanish phonology would predict, and it points to the consonants as
the place to look instead.

\subsection{/s/ in coda}
\label{sec:coda}

If dialectal variation in Spanish lies primarily in the consonants, the obvious
candidate for contrast within the Iberian Peninsula is the behaviour of /s/ in
syllable-final position. Descriptions of the Southern speech consistently point
to its weakening (aspiration, assimilation to the following segment, or elision)
in contrast to the full realization retained by Northern speech
\citep{narbona1998,gerfen2002,torreira2007}.

After Section~\ref{sec:pseudo}, this hypothesis has one feature that makes it
especially attractive: it is tested inside each speaker. The comparison that
matters runs between the /s/ in coda and the /s/ in onset of the same person,
rather than between a Northern /s/ and a Southern one. Each speaker acts as their
own control, so an imbalance in clips per speaker cannot manufacture the effect,
because it moves both terms of the comparison equally. Without phonemic
segmentation, the preliminary analysis had no way to ask this question.

The alignment yields 21,639 /s/ segments from 890 speakers, of which 12,269 are
coda segments and 9,370 are onset segments. For each segment, the spectral center
of gravity (CoG) is measured along with duration and root-mean-square energy.

\begin{table}[t]
\centering
\small
\setlength{\tabcolsep}{4pt}
\begin{tabular}{llrrrr}
\toprule
Position & Measurement & North & South & $d$ & \makecell{$p$\\(Holm)} \\
\midrule
Coda  & CoG (Hz) & 3,118.7 & 2,417.1 & $+0.56$ & $<$0.001 \\
Coda  & Duration & --      & --      & $+0.30$ & 0.019 \\
Coda  & RMS      & --      & --      & $+0.07$ & 1.000 \\
Onset & CoG (Hz) & 3,517.0 & 3,450.1 & $+0.06$ & 1.000 \\
Onset & Duration & --      & --      & $-0.27$ & 0.033 \\
Onset & RMS      & --      & --      & $+0.03$ & 1.000 \\
\bottomrule
\end{tabular}
\caption{/s/ in North and South, per syllabic position, in speaker level. The
dialectal contrast appears in the coda and disappears in the onset.}
\label{tab:coda}
\end{table}

The pattern matches the hypothesis rather than a general difference in how these
speakers produce sibilants. In coda, the Southern CoG sits 700 Hz below the
Northern one, with $d = +0.56$. In onset, the same varieties differ by 67 Hz,
with $d = +0.06$ and $p = 1.000$. The interaction between dialect and position on
the CoG is significant ($+196.9$ Hz, $p = 0.006$): the South does not produce
different sibilants, it produces different codas, which is what a phonological
process conditioned by position looks like.

The contrast can be summarized in a single measure per speaker, the weakening
index, defined as the difference between the CoG of their onsets and that of
their codas: the higher the value, the more the coda is weakened relative to the
full realization by the same speaker. The three varieties are ranked as follows:
North, $+359$ Hz over 369 speakers; Mexico, $+836$ Hz over 249; and South,
$+1,022$ Hz over 165.

When comparing the South versus the North and Mexico combined, $d = +0.63$ with
$p < 0.0001$: the strongest effect in the entire study, and the only finding that
dialectology would have predicted in advance. The ranking also replicates what
has been described in the literature without having been imposed by the study
design: Northern peninsular Spanish is the most conservative variety in coda
position, Southern peninsular Spanish is the most innovative, and Mexican Spanish
occupies an intermediate position, consistent with a distribution in which
coastal varieties aspirate and highland varieties maintain
\citep{hualde2005,arias2007}.

\begin{figure*}[t]
\centering
\includegraphics[width=\textwidth]{fig_codas.png}
\caption{Weakening of /s/ in coda: Spectral center of gravity and duration of /s/
according to syllabic position in the three varieties, and F1 of the preceding
vowel according to context. The distinction between varieties appears only in the
coda.}
\label{fig:codas}
\end{figure*}

This is the positive result of the study, and I return to what it implies
methodologically in Section~\ref{sec:discussion}.

\subsection{The vowel before coda /s/}
\label{sec:precoda}

The compensatory accounts introduced in Section~\ref{sec:dialectal} can now be
tested directly, by contrasting within each speaker the vowels that precede a
coda /s/ with the same vowels in other contexts. A total of 815 speakers provided
data in both contexts: 397 from the North, 168 from the South, and 250 from
Mexico.

The interaction between dialect and context is significant for F1 ($+0.101$,
$p = 0.005$), for F2 ($+0.070$, $p = 0.036$), and for duration ($+0.003$ s,
$p = 0.036$). There is, therefore, a vowel effect associated with the /s/ coda
context that differs across varieties.

Its interpretation, however, requires precision, because it does not match what
is described in the literature. The intra-speaker shift in F1 when moving to a
/s/ coda context is $-0.146$ in the North and $-0.086$ in the South, in Lobanov
z-units. Both values are negative: both varieties close the vowel before the /s/
coda, and the South simply closes it less. There is no Southern vowel opening in
these data; there is a generalized closing in which the South participates to a
lesser extent. The effect is also small, $d = -0.17$, and does not reach
significance on its own ($p = 0.063$). Nor is compensatory lengthening observed:
the difference in duration goes in the same direction but with a smaller
magnitude ($d = -0.10$, $p = 0.241$).

This result changed a lot with sample size. With 34 Southern speakers the same
contrast gave $d = -0.83$, a large effect that fit the traditional description
nicely. With 168 speakers it is $-0.17$. The interaction is still significant,
but the magnitude has dropped to a fifth.
This is a textbook case of regression to the mean: small samples do not produce
biased effects in a specific direction; rather, they produce extreme effects, and
those that survive the filter of attention are systematically the largest. It is
the same phenomenon that Section~\ref{sec:pseudo} documents from another angle,
and it is worth keeping this in mind when reading any phonetic study based on
just a few dozen speakers.

\subsection{Exploratory analyses}

Two unsupervised analyses round out the picture and, incidentally, serve to
correct certain conclusions from the preliminary analysis.

By applying $k$-means as implemented in scikit-learn \citep{pedregosa2011} to the
normalized formants and evaluating the resulting partition against the actual
vowel labels using the adjusted Rand index \citep{hubert1985}, normalization
improves the recovery of the vowel categories (ARI = 0.424 with Lobanov versus
0.332 in raw hertz) while the inclusion of bandwidths degrades it (0.357 and
0.258, respectively). Bandwidths do not contribute to vowel structure, they
dilute it.

The preliminary analysis had concluded, based on the silhouette coefficient, that
Spanish is organized into ``three perceptual zones'' rather than five vowels. The
$k$-scan reproduces that result, since the silhouette does peak at $k = 3$
(0.454), but the ARI places it at $k = 4$ (0.479), and the actual number of
categories is five. The discrepancy has a technical, rather than phonological, explanation:
the silhouette measure favors partitions with few compact clusters and decreases
mechanically as dimensions are added, so it is not comparable across spaces of
different dimensionality; the ARI, being measured against a known truth, is
comparable \citep{rousseeuw1987}. The $k = 3$ result was a property of the metric,
not of Spanish.

When training four classifiers to predict dialect based on acoustic parameters,
using speaker-disjoint cross-validation folds, the balanced accuracy is 0.500 in
all four cases, that is, exactly random. Contrasts that include Mexico are only
slightly higher (0.527 and 0.530).

The positive control shows that this comes from the problem and not from the
features: the same models, on the same data and the same partitions, predict the
vowel identity of speakers they have never heard with an accuracy of 0.762
against a chance level of 0.200, so 3.8 times chance. The measures carry plenty
of phonetic information. They carry no dialectal information.

The experiment also provides a second demonstration of the argument in
Section~\ref{sec:pseudo}. If the partitions are made randomly across tokens,
allowing the same speaker to appear in both the training and test sets, the
accuracy in dialect prediction is inflated by $+0.20$ on a sample constructed
like the one in the preliminary analysis. On the current sample, the same
comparison yields an increase of $+0.004$. Speaker identity leakage only distorts
the results when a few speakers dominate the sample. When the design is balanced,
both validation schemes converge.

\begin{figure*}[t]
\centering
\includegraphics[width=\textwidth]{fig_classification.png}
\caption{Dialect prediction against vowel identity prediction with
speaker-disjoint cross-validation and with random partitions. Dialect sits at
chance under both schemes in the balanced sample. Vowel identity is recovered far
above chance for previously unseen speakers.}
\label{fig:classification}
\end{figure*}

% =============================================================================
\section{Discussion}
\label{sec:discussion}

In the vowel space, Northern and Southern Peninsular Spanish are
indistinguishable at speaker level. In the treatment of coda /s/ they separate
clearly, with the onset serving as a null control. The large effects that the
preliminary analysis had located in the vowels did not exist at all.

One positive result and two null ones is not a coincidence, and the reason lies
in the design rather than in dialectology. The contrast that survives is the only
one drawn inside the speaker. Comparing the /s/ in coda with the /s/ in onset of
the same person neutralizes everything in this corpus that travels with dialect:
vocal tract anatomy, microphone, room noise, age, gender and reading style.
Comparisons between groups have to handle all of that through normalization,
statistical control and balanced sampling. In a corpus of volunteers, a
well-chosen intra-speaker contrast is worth more than any statistical refinement
applied to an overall comparison.

This also sets out what Common Voice is good for in phonetics. The unevenness
described in Section~\ref{sec:corpus} is not a defect to be tolerated: it governs
how any study on this corpus has to be designed, which makes the speaker
identifier the variable that decides whether an analysis is valid. A study that
samples clips without grouping them by speaker risks measuring individuals while
believing it measures varieties, and the risk grows with the corpus, because more
data buys significance without buying independence. With per-speaker caps and
aggregation beforehand, the same corpus supports solid inference.

One caveat remains. Section~\ref{sec:vowelspace} showed that no difference is
detectable, which is not the same as showing that none exists, and the gap
between those two statements matters here. The descriptive literature does not
call the Southern vowel system slightly different from the Northern one; it calls
it audibly and systematically different. That stronger claim is the one these
data fail to support, both in the vowel space as a whole and in the compensatory
opening before coda /s/, where the effect is real but five times smaller than a
small sample had suggested. What the data do support is the consonantal claim,
which carries the largest effect in the study.

% =============================================================================
\section{Limitations and conclusions}
\label{sec:limitations}

The data consists of read speech produced by volunteers under uncontrolled
recording conditions; therefore, the conclusions do not automatically apply to
spontaneous speech. Dialect labels are self-reported and unverifiable, and their
granularity is coarse: the label ``Southern Peninsula'' encompasses Andalusian,
Extremaduran, and Murcian varieties whose phonetic behavior is not homogeneous,
which necessarily attenuates any contrast. The sex imbalance is mitigated through
normalization but does not disappear.

On measurement, the validation with synthetic vowels shows a bias of +33.7\% in
the F1 estimated for /u/, so nothing should be concluded from that vowel; the
bias is identical in both dialects and cannot create spurious differences, but it
does add noise. In the /s/ analysis, forced alignment marks a boundary for the
segment even when it has been elided completely, so a full elision is measured as
a segment with a very low center of gravity rather than as an absence, which
probably understates the Southern weakening rather than exaggerating it. And the
null result for the vowel space constrains without excluding: differences below
$d \approx 0.24$ are still compatible with these data.

Three conclusions, in the order the study reached them.

First, the large effects in the initial analysis (F1 at $d = 1.34$ and B2 at
$d = -1.43$) came from the sampling rather than from the dialects. Two
speakers supplied 54\% and 80\% of their respective cells, and the size of the
effects tracks how much those two are allowed to contribute. Once the speaker
becomes the unit of inference, nothing is left of them.

Second, across 461 Northern and 186 Southern speakers I find no detectable
difference in the vowel space, in a design that would have caught anything at
$d \geq 0.24$. The Mexican outgroup rules out the obvious alternative
explanation: the same procedure, applied to a contrast nobody disputes, produces
effects 2.5 times larger.

Third, what separates these two varieties is consonantal. Southern speakers
weaken coda /s/ substantially and consistently ($d = +0.63$), onset /s/ shows
nothing, and the three varieties line up as the dialectological literature would
predict. This was also the only contrast I tested within speakers rather than
between groups, and I doubt that is a coincidence. On a corpus like this one, a
comparison that each speaker controls for themselves is worth more than a more
elaborate analysis of one that nobody controls.

% =============================================================================
\section*{Declaration on the use of LLMs}

A large language model was used as a writing aid for this report, restricted to
language editing: grammar, spelling, and phrasing of text I had already written,
together with the correction of errors introduced when I translated my own draft
from Spanish into English. It was also used to flag repeated passages so that I
could shorten them.

The research question, the study design, the sampling and analysis code, the
statistical decisions, the results and their interpretation are my own work. No
text was generated from a prompt and inserted without my having written the
corresponding content first, and every number reported here was produced by my
own pipeline and checked against its output tables.

I keep a full record of the editing sessions and can provide it on request.

% =============================================================================
\bibliography{references}
\bibliographystyle{acl_natbib}

\end{document}
